\chapter{Experimental Results}

\vspace*{1cm}

This chapter presents the experimental results obtained after running the
fine-tuning benchmark on a machine equipped with a performant GPU. The objective
is to compare the optimizers according to both model quality and computational
efficiency. The benchmark was executed with a small experimental setup of 10
training steps and a batch size of 4, which makes it possible to compare the
behavior of the optimization methods under the same conditions.

\section{Benchmark Setup}

The benchmark compares several optimization configurations applied to
\texttt{google/flan-t5-small}. In addition to the main optimizers studied in the
project, the results also include an inference-only baseline and an AdamW variant
with gradient clipping and learning-rate scheduling.

\vspace*{0.5cm}

The compared configurations are:

\begin{itemize}
    \item Inference Only baseline.
    \item SGD.
    \item SGD with momentum.
    \item Adam.
    \item AdamW with gradient clipping and scheduler.
    \item L-BFGS.
    \item Newton-CG.
\end{itemize}

\section{Global Comparison}

Figure~\ref{fig:optimizer_comparison_matrix} summarizes the main results of the
optimizer benchmark. It includes the training loss convergence, validation loss,
training step latency, peak VRAM usage, training throughput, and inference
latency.

\begin{figure}[H]
    \centering
    \includegraphics[width=\textwidth]{chapters/images/optimizer_comparison_matrix.png}
    \caption{Comparison matrix of the optimizer benchmark results.}
    \label{fig:optimizer_comparison_matrix}
\end{figure}

\section{Numerical Results}

The numerical results are summarized in Table~\ref{tab:benchmark_results}. Lower
values are better for loss, step latency, inference latency, VRAM usage, and RAM
increase. Higher values are better for throughput.

\begin{table}[H]
    \centering
    \scriptsize
    \begin{tabular}{lrrrrrr}
        \toprule
        \textbf{Optimizer} & \textbf{Train Loss} & \textbf{Val Loss} &
        \textbf{Step ms} & \textbf{Inf. ms} & \textbf{Tok/s} &
        \textbf{VRAM MB} \\
        \midrule
        Inference Only & 10.9235 & 10.8968 & 76.43 & 224.15 & 20096.52 & 553.23 \\
        SGD & 10.8635 & 10.8049 & 84.97 & 211.06 & 18076.50 & 1459.34 \\
        SGD + Momentum & 10.7055 & 10.5485 & 92.56 & 216.98 & 16595.08 & 1714.17 \\
        Adam & 9.4310 & 9.2548 & 190.80 & 1246.35 & 8050.34 & 1969.00 \\
        AdamW + Clip + Sched & 10.2606 & 10.0681 & 112.17 & 217.35 & 13693.65 & 1969.00 \\
        L-BFGS & 10.9365 & 10.8968 & 248.46 & 228.55 & 6182.03 & 5281.85 \\
        Newton-CG & 10.9030 & 10.1138 & 1445.42 & 1276.49 & 1062.67 & 6334.53 \\
        \bottomrule
    \end{tabular}
    \caption{Numerical summary of the optimizer benchmark.}
    \label{tab:benchmark_results}
\end{table}

\section{Training Loss Analysis}

The training loss curve shows that Adam achieves the strongest decrease during
the benchmark. Its final training loss is 9.4310, which is clearly lower than the
other optimizers. This indicates that Adam adapts the model parameters more
effectively during the limited number of training steps.

\vspace*{0.5cm}

SGD with momentum also improves the loss compared with plain SGD. Its final
training loss is 10.7055, while SGD reaches 10.8635. This confirms the advantage
of adding momentum to stabilize and accelerate gradient-based updates.

\vspace*{0.5cm}

L-BFGS and Newton-CG do not provide a clear improvement in this experiment. Even
though they are theoretically more advanced optimization methods, their practical
cost is high for transformer fine-tuning and their loss values remain close to
the baseline in this short benchmark.

\section{Validation Loss Analysis}

The validation loss confirms the same general trend. Adam obtains the best
validation loss with a value of 9.2548. This means that, among the tested
optimizers, Adam provides the best fine-tuning quality on the validation samples.

\vspace*{0.5cm}

AdamW with clipping and scheduler gives the second-best validation loss, with
10.0681. Newton-CG also reaches 10.1138, but this result must be interpreted
carefully because it requires a much higher computational cost. SGD with momentum
is better than plain SGD, while L-BFGS remains very close to the inference-only
baseline.

\section{Latency and Throughput Analysis}

From an efficiency point of view, Adam is not the fastest optimizer. Its average
training step latency is 190.80 ms, which is higher than SGD, SGD with momentum,
and AdamW with clipping and scheduler. Adam also has a much higher inference
latency in the benchmark, reaching 1246.35 ms per batch.

\vspace*{0.5cm}

SGD provides the best practical latency among the trained optimizers, with 84.97
ms per training step and 211.06 ms inference latency. It also keeps a high
throughput of 18076.50 tokens per second. SGD with momentum remains efficient,
with 92.56 ms per step and 16595.08 tokens per second.

\vspace*{0.5cm}

Newton-CG is the slowest method in the benchmark. Its average step latency is
1445.42 ms and its throughput is only 1062.67 tokens per second. This confirms
that second-order methods are expensive for transformer fine-tuning.

\section{Memory Usage Analysis}

The peak VRAM usage also shows important differences between optimizers. The
inference-only baseline uses 553.23 MB of VRAM. Among the trained optimizers, SGD
uses the least GPU memory with 1459.34 MB. SGD with momentum uses 1714.17 MB, and
Adam uses 1969.00 MB.

\vspace*{0.5cm}

L-BFGS and Newton-CG require much more GPU memory. L-BFGS reaches 5281.85 MB,
while Newton-CG reaches 6334.53 MB. These values make them less suitable for Edge
AI contexts, where memory is usually limited.

\section{Discussion of the Best Compromise}

The results show that Adam gives the best model quality because it achieves the
lowest training and validation losses. However, this comes with higher latency and
lower throughput compared with SGD-based methods. Therefore, Adam is the best
choice when the priority is fine-tuning quality.

\vspace*{0.5cm}

For an Edge AI perspective, SGD and SGD with momentum are more efficient. SGD is
faster and uses less VRAM, while SGD with momentum provides a better validation
loss than plain SGD with only a moderate increase in latency and memory usage.
AdamW with clipping and scheduler also gives an interesting balance between
quality and efficiency.

\vspace*{0.5cm}

L-BFGS and Newton-CG are less practical in this experiment. Their computational
and memory costs are high, and their performance does not justify this overhead
for the current fine-tuning setup.

\section{Conclusion}

The benchmark highlights a clear tradeoff between quality and efficiency. Adam is
the most effective optimizer in terms of loss reduction, but SGD-based methods
are more efficient for latency, throughput, and memory usage. For future
compression and Edge AI deployment, SGD with momentum or AdamW with clipping and
scheduling can be considered as good compromise candidates, while Adam remains
the strongest quality baseline.
